%======================================================================
\section{Introduction}\label{sec:intro}
%======================================================================

Many financial derivatives combine path dependence with early-exercise features. For instance, an Amerasian put pays the strike minus the running arithmetic average upon exercise, and a similar structure arises in Bermudan options on averages, lookback options, and callable structured notes. At each exercise date, the holder compares the immediate exercise payoff with the continuation value, both of which depend on the price history; pricing thus reduces to an optimal stopping problem on path space. Classical lattice and PDE methods require a finite-dimensional state variable that summarizes this history. However, augmenting the state with running averages or extrema must be tailored to each payoff and increases the effective dimension, a difficulty shared by regression-based Monte Carlo methods \cite{longstaff2001} through the choice of basis functions.

Deep backward stochastic differential equation (BSDE) methods eliminate the spatial grid. The method of Han, Jentzen, and E \cite{han2018} simulates the forward process, represents the hedge integrand $Z$ by a neural network, and minimizes the terminal mismatch; it scales to high dimensions and returns a hedge together with the price. Early exercise is usually handled by backward schemes \cite{hure2020}, which fit one regression per exercise date, proceed from maturity to the initial date, and freeze each fitted layer before training the previous one. Bayraktar, Feng, and Zhang \cite{bayraktar2024} add a signature layer to this backward scheme and price path-dependent American options with a convergence rate in the truncation order. However, backward schemes require $K$ sequential training stages for $K$ exercise dates, so the approximation error of each frozen layer propagates into the targets of all preceding layers.

Huang and Oosterlee \cite{huang2026} proposed the compound BSDE as a fully forward alternative. The horizon is split at the exercise dates into $K$ layers, each carrying its own BSDE, and consecutive layers are linked by the interface condition $Y_{j,T_j}=\max\{\Phi_j,Y_{j+1,T_j}\}$. A single loss penalizes all interface residuals and the terminal residual, and one optimizer updates all layers simultaneously. Nevertheless, three aspects of this construction limit its scope. First, the networks take only the current state as input, which excludes path-dependent payoffs. Second, the loss is differentiated through both sides of every interface. Because the interface residual is a Bellman residual, minimizing it jointly over both arguments yields stationary points that generally differ from the solution of the Bellman equation \cite{baird1995,sutton2018}; whether this residual-gradient bias affects the compound BSDE has not yet been examined. Third, its a posteriori error estimate requires a smallness condition on the Lipschitz constant of the driver.

This paper develops a compound BSDE for path-dependent early-exercise contracts and studies how its interfaces should be trained. Local log-signatures of the time-augmented path drive a neural rough differential equation (Neural RDE) \cite{morrill2021}, whose hidden state summarizes the history and serves as input to the value and hedge heads of every layer. A single encoder is shared by all layers, and its hidden state is not reset at the exercise dates. The interface target $\max\{\Phi_j,\hY_{j+1,T_j}\}$ enters the loss through a stop-gradient operator, so each value head regresses on a fixed target, as in fitted value iteration and target networks \cite{mnih2015}. All layers are still updated in every optimizer step, and training remains a single forward pass. Our main contributions are threefold.

\emph{Stationary points of interface training.} Under stop-gradient training, every value head is the conditional expectation of its target given the encoder filtration (\cref{thm:stationary}). The stationary price therefore coincides with the optimal stopping value $V_0^{(M)}$ under this filtration, and it equals the price of the monitored problem whenever the log-signature captures the continuation values. By contrast, joint differentiation introduces a bias term proportional to the residual of the preceding layer on the predicted continuation region, and a two-period example shows that the resulting price error can be negative. In our experiments with $K=4$ to $24$ exercise dates, joint training overprices Bermudan and Amerasian puts by $1.4$--$9.5\pct$; this bias shrinks only as the learning step is refined, not as the training budget increases. Stop-gradient training stays within $0.32\pct$ for four exercise dates and near $1\pct$ for up to $24$ dates.

\emph{Error analysis of the path-dependent compound BSDE.} The compound BSDE admits a unique solution whose layers are Lipschitz functionals of the path, with a constant independent of $K$ (\cref{thm:wellposed}). The training loss bounds the squared price and hedge errors of any adapted output up to a factor of $2K$ (\cref{thm:aposteriori}). Discrete discounting turns the forward Euler scheme into an exact martingale transform, which removes the smallness condition of \cite{huang2026}. Moreover, the total error splits into a learning-grid term, a monitoring term, a truncation term, and a network term (\cref{thm:attainable}). Level-two log-signatures capture all information required by arithmetic Asian and Amerasian payoffs, and level $m=h/\hf$ recovers the entire monitored path (\cref{lem:information}).

\emph{Certified prices.} From a single forward pass, the trained network yields both an exercise rule and a dual martingale. The resulting interval $[L_0,U_0]$ contains the price, and its width is at most $\sqrt{K\cL(\theta)}+(K-1)\eex$ (\cref{thm:primaldual}). On Bermudan, Amerasian, basket, and max-call options, the certified intervals match or improve on Longstaff--Schwartz regression \cite{longstaff2001}, deep optimal stopping \cite{becker2019}, and signature-based primal--dual stopping \cite{bayer2025primal}, and the accuracy remains nearly constant from one to ten assets. Under rough Bergomi dynamics, where the price is no longer a Markov state, the log-signature encoder yields tighter dual bounds than both a Markov-state encoder and signature regression.

Our encoder builds on signature methods for path-dependent problems. Signatures provide universal coordinates for path functionals \cite{lyons1998,cuchiero2025}. They have been combined with deep BSDE solvers in forward \cite{feng2023,sun2025} and backward \cite{bayraktar2024} form, as well as with regression and duality for non-Markovian stopping \cite{bayer2025primal,bayer2026stopping}. Neural controlled and rough differential equations evolve a hidden state driven by the path or by its local log-signatures \cite{kidger2020,morrill2021}. Fang, Ni, and Wu \cite{fang2026} solve path-dependent parabolic PDEs with Neural RDEs, and Alzahrani \cite{alzahrani2025} attaches BSDE heads to a log-signature Neural RDE. These works, however, treat a single terminal condition. The present work integrates one shared log-signature Neural RDE encoder into a multi-layer optimal stopping architecture and analyzes the interplay between the path representation and interface training.

The paper is organized as follows. \Cref{sec:setting} introduces the compound BSDE on path space. \Cref{sec:method} presents the encoder, the forward recursion, and the stop-gradient loss. \Cref{sec:analysis} contains the error analysis, and \cref{sec:experiments} reports the numerical experiments. \Cref{sec:conclusion} concludes, and full proofs are collected in \cref{app:proofs}.
